\section*{ЗАКЛЮЧЕНИЕ}
\addcontentsline{toc}{section}{ЗАКЛЮЧЕНИЕ}

В результате выполнения выпускной квалификационной работы была разработана программно-информационная поисковая система с графическим интерфейсом, состоящая из библиотеки регулярных выражений на языке Си и десктопного поискового приложения на C++ с использованием фреймворка Qt.

Разработанная библиотека регулярных выражений представляет собой полностью функциональный однофайловый модуль на языке Си с переносимым C API, поддерживающий широкий синтаксис шаблонов: метасимволы, символьные классы, группы захвата, обратные ссылки, альтернацию, ленивые и жадные квантификаторы, модификаторы регулярных выражений, в том числе регистронезависимый поиск с поддержкой кириллицы в кодировке UTF-8. Поисковое приложение предоставляет пользователю текстовый редактор с подсветкой совпадений, интерактивный диалог поиска и замены, а также модуль рекурсивного поиска по файлам директории с отображением результатов и навигацией к найденным фрагментам.

Основные результаты работы:
\begin{enumerate}
	\item Изучены методы реализации обработчиков регулярных выражений (интерпретация синтаксического дерева, симуляция конечного автомата, рекурсивный спуск с возвратами), язык Си и принципы проектирования C API, фреймворк Qt.
	\item Разработана концептуальная модель библиотеки регулярных выражений. Определены требования к библиотеке, её C API и поисковому приложению.
	\item Спроектирована программная система библиотеки регулярных выражений: определены внутренние структуры данных, формальная грамматика поддерживаемого синтаксиса шаблонов, алгоритмы компиляции и сопоставления.
	\item Спроектирован переносимый C API библиотеки, обеспечивающий её интеграцию в программы на Си-подобных языках, в том числе в C++-проекты через блок \verb`extern "C"`.
	\item Реализована библиотека регулярных выражений средствами языка Си. Разработан встроенный набор автоматизированных тестов, охватывающий все семейства экспортируемых функций.
	\item Спроектирована программная система поискового приложения: определена архитектура модулей, схема интеграции библиотеки, механизмы взаимодействия компонентов.
	\item Реализовано поисковое приложение на языке C++ с использованием фреймворка Qt, включающее текстовый редактор, диалог поиска и замены и модуль рекурсивного поиска по директории.
\end{enumerate}

Все требования, объявленные в техническом задании, были полностью реализованы, все задачи, поставленные в начале разработки проекта, были также решены.

Таким образом, цель работы достигнута.